\section{Structured infinite fibres}\label{sec:8}
Sections 2–7 constrain the fibres that are \emph{finite}. This section describes the infinite ones. The result is not merely that an unbounded fibre exists, but that it contains an infinite subset generated by arithmetic data of one of three forms produced by the proof, and that the data can be exhibited. The certificate describes a subset, never a parametrisation of the whole fibre. The three forms are those the argument emits; they are not claimed to be canonical, and a fibre may carry certificates of more than one form.

For $P ∈ \mathrm{Poly} (n+1) r$ and $a ∈ ℤ^{n}$ let $\mathrm{fiberZZ} P a := \{c | ∃ u, P(a, c, u) = 0\}$ be the scalar fibre over a.

\begin{definition}[generic certificate; \texttt{InfiniteFiberCertificate.\allowbreak{}lean}]\label[definition]{definition:8.1}
An \emph{infinite-subset certificate} for (P, a) consists of an infinite index type, an injective output map from it into ℤ, and a witness map, together with proofs that each output lies in the fibre with the given witness. It certifies that the fibre is infinite by exhibiting it as containing the image of an injection.

\end{definition}
The word \emph{constructive} is deliberately avoided here, and the following four statements are kept apart throughout the paper.

\begin{longtable}{@{}>{\raggedright\arraybackslash}p{0.553\linewidth}>{\raggedright\arraybackslash}p{0.347\linewidth}@{}}
\toprule
statement & status \\
\midrule\endhead
the output family has the stated arithmetic form (progression, orbit, step orbit) & proved, with the data displayed \\
every output of the family lies in the fibre, i.\allowbreak{}e. a witness exists for it & proved \\
a witness map exists, assigning to each output one realising witness & proved, but the assignment is made extensionally (\texttt{Classical.\allowbreak{}choice}) from the membership proofs \\
the witness map is computable, or a certificate can be extracted algorithmically from P & not proved, and not claimed \\
\bottomrule
\end{longtable}
The certificate presents the infinite output subset by explicit arithmetic data. Witnesses realising individual outputs exist and are selected extensionally; no theorem here asserts an effective witness-generating algorithm.

\begin{definition}[the three kinds; \texttt{StructuredCertificates.\allowbreak{}lean}]\label[definition]{definition:8.2}
A \emph{structured certificate} \texttt{StructuredCertificate P a} is one of three realizations, each of which carries arithmetic data that \emph{pin} the output map rather than merely asserting that some injection exists.

\begin{enumerate}
\item \emph{Polynomial progression.} A polynomial $q ∈ ℤ[t]$ with $1 ≤ \mathrm{deg} q ≤ 2, a$ threshold t₀, witnesses $w_{t}$, and proofs that $P(a, q(t), w_{t}) = 0$ for every $t ≥ t_{0}$ and that q is injective on that tail. The fibre contains an arithmetic or quadratic progression outright.
\item \emph{Pell-unit orbit.} A positive nonsquare D, a unit (x₀, y₀) of x² − Dy² = 1 with $y_{0} ≠ 0, a$ base point (X₀, Y₀) of norm $N_{0} := X_{0}^{2} − DY_{0}^{2}$, and the orbit generated by the unit as an explicit recurrence,

\[
X_{k+1} = x_{0}X_{k} + D y_{0} Y_{k},  Y_{k+1} = y_{0}X_{k} + x_{0}Y_{k},
\]
together with an affine projection with denominator, $\text{denom}·\text{output}_{k} = p_{c} + p_{x} X_{k} + p_{y} Y_{k}$ with $\mathrm{denom} ≠ 0$, witnesses, and injectivity of the output. The norm is preserved along the orbit as a theorem of the data (\texttt{PellUnitOrbitRealization.\allowbreak{}orbit\_\allowbreak{}norm}), not as an extra hypothesis.
\item \emph{Moving-modulus step orbit.} Integers $b > 0, μ$, and a positive \emph{modulus factor} F, a sign ε ∈ \{±1\}, a sequence $c_{k}$ with moduli $t_{k} := b·c_{k} + μ > 0$, and primes $p_{k}$ with

\[
F·b·t_{k} < p_{k}, \quad p_{k} ≡ 1 (mod F·b·t_{k}), \quad t_{k+1} = t_{k} · p_{k},
\]
with witnesses for the outputs $ε·c_{k}$. Strict monotonicity of $c_{k}$ is a consequence of the recurrence and does not have to be assumed (\texttt{MovingModulusStepOrbitRealization.\allowbreak{}c\_\allowbreak{}strictMono}).
\end{enumerate}
\end{definition}

The three kinds are a sufficient certificate language, not a classification: no claim is made that they are canonical, exhaustive of the fibre, or mutually exclusive, and a fibre may carry certificates of more than one kind (\cref{example:8.6} exhibits two). The third kind deserves a word. The relevant family is an iterated moving-modulus orbit rather than a fixed prime progression: at each step the congruence that the next prime must satisfy is taken modulo a quantity that itself grows with the current output, so the object is an orbit under a step map whose modulus moves. Dirichlet's theorem is what makes each individual step possible — there is always a prime in the relevant progression — but the family is generated by iterating the step from a base point, not by ranging over the primes of one progression. Writing the certificate as data with the recurrence \texttt{t\_\allowbreak{}\{k+1\} = t\_\allowbreak{}k·p\_\allowbreak{}k} in it makes the distinction unmissable.

\begin{theorem}[soundness; \texttt{fiber\_\allowbreak{}infinite\_\allowbreak{}of\_\allowbreak{}structuredCertificate}, \texttt{StructuredCertificate.\allowbreak{}toSubset}]\label[theorem]{theorem:8.3}
Every structured certificate yields a generic certificate and hence proves the fibre infinite; a certified fibre exceeds every bound (\texttt{fiber\_\allowbreak{}unbounded\_\allowbreak{}of\_\allowbreak{}certificate}); certificates pull back along integral changes of witness coordinates (\texttt{certificate\_\allowbreak{}pullback}), which is what allows a certificate produced in transformed coordinates to be reported in the original ones.

\end{theorem}
\begin{proof}
Three assertions, each by the data.

\emph{Every structured certificate yields a generic one.} A generic certificate asks for an infinite index type, an injective output map into ℤ, a witness map, and proofs that each output lies in the fibre with its witness. Each kind of \cref{definition:8.2} supplies exactly these. For a polynomial progression the index type is the tail \{t : t ≥ t₀\}, which is infinite; the output map is q, injective on that tail by hypothesis; the witnesses are the $w_{t}$, and $P(a, q(t), w_{t}) = 0$ is part of the data. For a Pell-unit orbit the index type is ℕ; the output map is determined by $\mathrm{denom}·output_{k} = p_{c} + p_{x} X_{k} + p_{y} Y_{k}$ with $\mathrm{denom} ≠ 0$, and its injectivity is part of the data; the witnesses are again given. For a moving-modulus step orbit the index type is ℕ and the outputs are $ε·c_{k}$, whose injectivity is not assumed but follows from the recurrence: $t_{k+1} = t_{k}·p_{k}$ with $p_{k} > F·b·t_{k} ≥ 1$ forces $t_{k}$, and hence $c_{k} = (t_{k} − μ)/b$, to increase strictly (\texttt{MovingModulusStepOrbitRealization.\allowbreak{}c\_\allowbreak{}strictMono}). In all three cases forgetting the arithmetic data and keeping the index type, the output map and the witnesses is a generic certificate.

\emph{A certified fibre is infinite, and exceeds every bound.} The output map of a generic certificate is an injection of an infinite type into the fibre, so the fibre contains an infinite set of integers. An infinite set of integers is unbounded in absolute value: a bound B confines it to the finite set [−B, B]. Hence no B bounds the fibre.

\emph{Pull-back along integral changes of witness coordinates.} Let T be an integral matrix with integral inverse and let P′(a, c, u) := P(a, c, Tu). Then u is a witness for P′ at output c precisely when Tu is one for P, and the correspondence u ↦ Tu is a bijection of ℤʳ. Given a certificate for P′, replace each witness w by Tw and keep the index type, the output map and the arithmetic data unchanged; every membership proof transports along the identity P′(a, c, w) = P(a, c, Tw). Injectivity of the output map is untouched because outputs are not changed. This is what allows a certificate produced in transformed coordinates — the coordinates in which the projection dichotomy works — to be reported in the original ones.

\end{proof}
\begin{definition}[the summit target; \texttt{SummitTarget.\allowbreak{}lean}]\label[definition]{definition:8.4}
For $P$ of total degree at most two,
\[
\begin{gathered}
\mathrm{UniformQuadraticFiberStructureZZ}\ P :⟺ ∃ C ∈ ℕ, ∀ a ∈ ℤ^{n}, \\
( ∀ c ∈ \mathrm{fiberZZ}\ P\ a, |c| ≤ C·(1 + ‖a‖_{∞})^{2} ) \ ∨ \ \mathrm{Nonempty}\ (\mathrm{StructuredCertificate}\ P\ a).
\end{gathered}
\]
\end{definition}

The order of the quantifiers is the content of the definition: C is chosen after P and \emph{before} a, so one constant serves every parameter point. The weak form ∀ a ∃ $C_{a}$, in which the bound may be chosen after seeing the parameter, is strictly weaker and is not what is proved here. The two horns cannot both hold at a given a, since a certificate exhibits an unbounded output family inside the fibre while the first horn bounds every element of it (\texttt{uniform\_\allowbreak{}bound\_\allowbreak{}certificate\_\allowbreak{}exclusive}). The weaker per-point form, in which the bound B may depend on a, is \texttt{QuadraticFiberBoundedOrStructured}.

The distinction matters and is not cosmetic. On a \emph{finite} fibre the per-point form is satisfied by taking B to be the fibre's own maximum, whatever that maximum happens to be (\texttt{weak\_\allowbreak{}target\_\allowbreak{}trivial\_\allowbreak{}on\_\allowbreak{}finite\_\allowbreak{}fibre}), so a theorem stated in the weak form would carry no information about how the bound depends on the parameters. The uniform form asserts the quadratic dependence, which is exactly the content of Section 2.

\begin{theorem}[structured summit, unconditional; \texttt{SummitClosure.\allowbreak{}lean}, \texttt{uniformQuadraticFiberStructureZZ\_\allowbreak{}of\_\allowbreak{}quadratic}, \texttt{singleQuadratic\_\allowbreak{}fiber\_\allowbreak{}bounded\_\allowbreak{}or\_\allowbreak{}structurallyInfinite}]\label[theorem]{theorem:8.5}
For every $P ∈ \mathrm{Poly}\ (n+1)\ r$ with $\mathrm{totalDegree}\ P ≤ 2$, $\mathrm{UniformQuadraticFiberStructureZZ}\ P$ holds. Consequently every scalar fibre of every equation of total degree at most two is bounded by $C(1+‖a‖_{∞})^{2}$ with $C$ depending only on $P$, or carries a structured certificate of one of the three kinds of \cref{definition:8.2}.

\end{theorem}
\begin{proof}
The proof has four steps, each of which is a separate theorem in the formalization.

\emph{(i) The uniform constant is free.} For quadratic P the strong target is equivalent to the certificate supply on the infinite side alone: UniformQuadraticFiberStructureZZ P ⟺ (∀ a, fiberZZ P a infinite → the fibre carries a structured certificate) (\texttt{uniformQuadraticFiberStructureZZ\_\allowbreak{}iff\_\allowbreak{}infiniteFiberStructured}). One direction is immediate; for the other, take the constant C supplied by \cref{theorem:2.1} in the sup-norm form, and split on whether the fibre is finite. If it is, \cref{theorem:2.1} gives the bound with that C; if it is not, the hypothesis supplies the certificate. So nothing is lost by demanding uniformity, and the whole remaining obligation is the infinite side.

\emph{(ii) Reduction to the nonnegative halves.} An infinite subset of ℤ is infinite on the nonnegative or on the nonpositive side (\texttt{infinite\_\allowbreak{}integer\_\allowbreak{}set\_\allowbreak{}pos\_\allowbreak{}or\_\allowbreak{}neg}). Reflecting the output coordinate maps one side to the other, and the reflection acts on certificates kind-preservingly: the polynomial q is replaced by its reflection, the affine projection of a Pell orbit by its negative, and the moving-modulus sign field ε flips (\texttt{OutputNegation.\allowbreak{}lean}, \texttt{structuredCertificate\_\allowbreak{}of\_\allowbreak{}flipZero}). It therefore suffices to certify a half-fibre that is infinite, with the sign carried as a parameter ε satisfying $ε^{2} = 1$.

\emph{(iii) Homogenisation: a half-fibre is the fibre of a homogeneous form.} This is the same decomposition used in the proof of \cref{theorem:2.1}, now used to transport structure rather than a bound. Relabelling and splitting into homogeneous components writes the specialised equation as H(x) + ℓ(a)·x + q(a) with H homogeneous of degree two in $x = (c, u)$, and the $ε = −1$ case is handled by the reflected form (\texttt{flipZero H}, \texttt{fiberVal\_\allowbreak{}flip}). One obtains H, an integer vector L and an integer Q with

$\mathrm{fiberSetG} H L Q = \{m ∈ ℕ | ε·m ∈ \mathrm{fiberZZ} P a\}$

(\texttt{exists\_\allowbreak{}fiberSetG\_\allowbreak{}eq\_\allowbreak{}fiberNatSigned}). The left-hand side is exactly the object about which \cref{sec:9} speaks.

\emph{(iv) The dichotomy in data-carrying form, and the conversion.} \cref{theorem:9.3} below is the projection dichotomy strengthened so that the infinite horn produces not the proposition "this set is infinite" but a \emph{structured family} inside it — a progression, a quadratic progression, a Pell orbit with its recurrence, or a moving-modulus orbit with its primes (\texttt{StructuredFamily}, \texttt{quadraticProjectionDichotomyZ\_\allowbreak{}all\_\allowbreak{}data}). Transporting such a family along the set equality of step (iii) and converting it into a certificate is then mechanical: choice supplies the witnesses from the membership proofs, the sign ε travels into the polynomial coefficients, into the affine projection, or into the sign field according to the kind, and the derived side conditions of the Pell kind (D > 0, D not a square, y₀ ≠ 0) follow from the unit equation together with $x_{0} ≥ 2$ (\texttt{structuredCertificate\_\allowbreak{}of\_\allowbreak{}family}). If instead the bounded horn holds, the half-fibre is finite, contradicting the assumption of step (ii).

\end{proof}
\begin{example}[one instance of each kind, and an overlap]\label[example]{example:8.6}
Take $n = 0$, so that the fibre is a single set of integers.

\begin{enumerate}
\item \emph{Polynomial progression.} The equation $c − u^{2} = 0$ has fibre \{t² : t ∈ ℕ\}, a degree-two progression with output map $c_{t} = t^{2}$. This is also the example that refutes any attempt to describe every one-unknown fibre by unit orbits alone.
\item \emph{Pell-unit orbit.} The equation $c^{2} − 2u^{2} − 1 = 0$ has as its fibre the set of x-coordinates of $x^{2} − 2y^{2} = 1$. That fibre \emph{contains} the orbit of the unit 3 + 2√2 under the recurrence of \cref{definition:8.2}(2), taken from the base point $(X_{0}, Y_{0}) = (1, 0)$, with output projection the identity $(\mathrm{denom} = 1)$: the outputs are 1, 3, 17, 99, …. The fibre is strictly larger than this orbit — it also contains −1, −3, −17, … — which is the general situation: a certificate exhibits a structured infinite subset of the fibre, not a parametrisation of all of it.
\item \emph{Moving-modulus step orbit.} For the singular branch with positive slope the construction of \cref{sec:9} returns the recurrence $b·c_{k+1} + μ = (b·c_{k} + μ)·p_{k}$ with the primes $p_{k}$ chosen stagewise above the branch's modulus factor; the outputs grow strictly and are pinned by that recurrence rather than by a closed form. An instance of the data, with $r = 3$ unknowns, is the equation
\end{enumerate}
\[
u^{2} + v^{2} + (c + 1)·w = 0,
\]
taking $b = 1, μ = 1, ε = 1$ and modulus factor $F = 2$, so that $t_{k} = c_{k} + 1$. Start from $c_{0} = 0, t_{0} = 1$. The step requires a prime $p_{k}$ with $p_{k} > F·b·t_{k} = 2t_{k}$ and $p_{k} ≡ 1 (\mathrm{mod} 2t_{k})$, which Dirichlet's theorem supplies at every stage; the first two choices are $p_{0} = 3$, giving $t_{1} = 3$ and $c_{1} = 2$, and $p_{1} = 7$, giving $t_{2} = 21$ and $c_{2} = 20$. Each output $c_{k}$ is realised by the witness $(u, v, w) = (0, 0, 0)$. This instantiates the constructor of \cref{definition:8.2}(3) exactly. It should be read as an exhibition of the data shape and not as evidence that the shape is forced: the fibre of this equation is all of ℤ, so the same fibre also carries a certificate of kind (1).
\end{example}

\emph{Overlap.} The three kinds are not mutually exclusive, and two fibres above carry certificates of two kinds each.

The fibre \{t²\} of (1) also carries a certificate of kind (2). Write $(3 + 2√2)^{k} = x_{k} + y_{k}√2$, so that $x_{k}^{2} − 2y_{k}^{2} = 1$. Take for the orbit the \emph{square} of that unit, $(x_{0}, y_{0}) = (17, 12)$, which satisfies $17^{2} − 2·12^{2} = 1$, and the base point $(X_{0}, Y_{0}) = (1, 0)$. The orbit of \cref{definition:8.2}(2) is then $X_{k} + Y_{k}√2 = (17 + 12√2)^{k} = (3 + 2√2)^{2k}$, so that

\[
X_{k} = x_{k}^{2} + 2y_{k}^{2} = 2x_{k}^{2} − 1,
\]
using $x_{k}^{2} − 2y_{k}^{2} = 1$. Hence $2·x_{k}^{2} = X_{k} + 1$, which is an affine projection with denominator in the sense of \cref{definition:8.2}(2):

\[
\text{denom} = 2, \quad p_{c} = 1, \quad p_{x} = 1, \quad p_{y} = 0, \quad \text{output}_{k} = x_{k}^{2}.
\]
The outputs are 1, 9, 289, … , each realised in the fibre of $c − u^{2} = 0$ by the witness $u = x_{k}$, and they are strictly increasing, hence injective. So the same fibre carries both a polynomial progression and a Pell-unit orbit. (Writing the projection against the \emph{unsquared} orbit would not do: $c = x_{k}^{2}$ is quadratic in the orbit coordinate, and \cref{definition:8.2}(2) admits only affine projections. Passing to the squared unit is what makes the projection affine.)

The fibre of the equation in (3) is all of ℤ, and it carries both the moving-modulus certificate displayed there and the polynomial progression $q(t) = t$ of kind (1).

A certificate is a witness of structure, not a canonical invariant of the fibre.

\begin{remark}[the shapes a fibre can have, and where each lands in \cref{theorem:8.5}]\label[remark]{remark:8.7}
All five occur, and each is exhibited in the formalization (\texttt{SummitInstances.\allowbreak{}lean}).

\end{remark}
\begin{longtable}{@{}>{\raggedright\arraybackslash}p{0.328\linewidth}>{\raggedright\arraybackslash}p{0.243\linewidth}>{\raggedright\arraybackslash}p{0.329\linewidth}@{}}
\toprule
fibre & example (n = 0) & horn of \cref{theorem:8.5} \\
\midrule\endhead
empty & 1 = 0 & bounded, vacuously \\
finite, nonempty & c² − 4 = 0, fibre \{−2, 2\} & bounded \\
the output zero alone & c = 0 & bounded \\
infinite on the nonpositive side only & c + u² = 0, fibre \{−t²\} & certificate, after the ε = −1 transport \\
infinite on both sides & 0 = 0, fibre ℤ & certificate \\
\bottomrule
\end{longtable}
The fourth row is why the proof works with signed halves: the nonnegative half of that fibre is \{0\}, so a construction looking only at nonnegative outputs would wrongly report a bounded fibre. The two halves meet exactly in the output zero (\texttt{halves\_\allowbreak{}inter}), so the fibre is infinite if and only if some half is, and the sign ε is carried through the construction rather than handled by a separate argument.

\begin{remark}[which kind comes from where]\label[remark]{remark:8.8}
The kind of certificate is determined by the branch of the proof of \cref{theorem:9.3} that produces it. Exhaustiveness is established by the typed case analysis inside the production theorem itself (\texttt{quadraticProjectionDichotomyZ\_\allowbreak{}all\_\allowbreak{}data}, \texttt{TopDichotomyData}); \texttt{BranchExhaustion.\allowbreak{}lean} is a generated documentation layer that mirrors that analysis as a finite table and does not independently establish it. The correspondence is many-to-one in both directions, and it is worth stating it as the table has it rather than in a slogan. In the full-rank layer the direction families give Pell-unit orbits, the $α = 0$ leaves give polynomial progressions, and the semidefinite leaf is bounded. In the singular block the positive-slope leaves give moving-modulus step orbits and the remaining leaves give polynomial progressions. In the one-witness layer — reached both by the $r ≤ 1$ dispatch and by the uncoupled rank-one leaf, which pass through Theorem V.4 — the leaves give polynomial progressions except the $Δ > 0$ non-square case, which gives a Pell-unit orbit, and the definite case, which is bounded. So no layer emits only one kind and no kind comes from only one layer; in particular Pell-unit orbits arise in the one-witness layer as well as in the full-rank one, which is what Corollary V.5 records. No branch is left without a kind, and no kind is produced by a branch that cannot support its data.

\end{remark}
